\documentclass[11pt]{amsart}

\usepackage[T1]{fontenc}
\usepackage{lmodern}
\usepackage{microtype}
\usepackage{mathtools,amssymb,amsthm}
\usepackage[hidelinks]{hyperref}

\hypersetup{
  pdftitle={The Maximum Euler-Sombor Index at Diameter Three, and Both Extrema at Nine Vertices}
}

\newtheorem{theorem}{Theorem}
\newtheorem{lemma}[theorem]{Lemma}
\newtheorem{proposition}[theorem]{Proposition}
\theoremstyle{remark}
\newtheorem{remark}[theorem]{Remark}

\DeclareMathOperator{\diam}{diam}
\newcommand{\EU}{\mathit{EU}}
\newcommand{\Gnd}{\mathcal G}

\title[Euler--Sombor at diameter three]
{The Maximum Euler--Sombor Index at Diameter Three, and Both Extrema at Nine Vertices}

\date{}

\subjclass[2020]{05C09, 05C35, 05C12, 92E10}
\keywords{Euler--Sombor index, degree-based topological index, diameter,
extremal graph, theta graph, bicyclic graph}

\begin{document}

\begin{abstract}
The Euler--Sombor index of a graph $G$ is
\[
 \EU(G)=\sum_{xy\in E(G)}\sqrt{d(x)^2+d(x)d(y)+d(y)^2}.
\]
Khanra and Das asked for the extremal values and extremal graphs of $\EU$
over ``the class of all connected graphs with given diameter''; Sekar,
Balachandran, Elumalai and Liu reproduce the request as their ``Problem~1'',
and a 2026 survey by Ali, Gutman, R\'eti, Albalahi and Hamza poses it with the
order fixed as well, which is the reading taken here. We settle one column and
one cell, not the problem. For every $n\ge4$ the maximum of $\EU$ over the
connected graphs of order $n$ and diameter exactly $3$ equals
\[
 \frac{\sqrt3}{2}(n-2)^2(n-3)+\sqrt{n^2-3n+3}+(n-3)\sqrt{3n^2-15n+19},
\]
attained by exactly one graph up to isomorphism, namely $H(1,n-3)$;
the proof is a closed form and uses no computation.
We also settle the single cell $n=9$, $D=3$ in the other direction: over the
$148\,229$ isomorphism classes of that cell the minimum is
$8\sqrt3+6\sqrt{19}=40.0098\ldots$, attained by exactly one class, the
\emph{bicyclic} theta graph $\Theta(3,3,4)$, so that the minimiser there is
neither a tree nor unicyclic and the best tree in the cell is only third. That
half is settled by the exhaustive census shipped with this paper, not by hand.
Problem~1 quantifies over all $n$ and all $D$ and is \emph{not} resolved here;
\S\ref{sec:open} says exactly what remains, including the one paywalled paper
that could reduce our maximum half to a rediscovery.
\end{abstract}

\maketitle

\section{The problem and its source}

For a graph $G$ write $d(x)=d_G(x)$. The \emph{Euler--Sombor index} is
\begin{equation}\label{eq:EU}
 \EU(G)=\sum_{xy\in E(G)}\sqrt{d(x)^2+d(y)^2+d(x)d(y)},
\end{equation}
in the normalisation printed on journal page~761 of \cite{Sekar} and in the
abstract of \cite{Kizilirmak}. Khanra and Das close \cite{KhanraDas} with a
list of open problems whose item~(1) reads, on journal page~546, verbatim:

\begin{quote}
Here we propose the following open problems on Euler Sombor index of a graph:
Determine the extremal values and extremal graphs with respect to Euler Sombor
index over the following sets: (1) the class of all connected graphs with given
diameter; \ldots
\end{quote}

Sekar, Balachandran, Elumalai and Liu \cite{Sekar} reproduce item~(1) as a
numbered problem. On journal page~762, immediately after the sentence ``Very
recently, Khanra et al.\ [17] posed the following open problem:'', they print
verbatim:

\begin{quote}
\textbf{Problem 1.} Determine the extremal values and extremal graphs with
respect to Euler Sombor index over the class of all connected graphs with given
diameter.
\end{quote}

The label ``Problem~1'' is theirs; the problem is Khanra and Das's. A 2026
survey by Ali, Gutman, R\'eti, Albalahi and Hamza \cite{Survey} still poses it,
as Problem~3.1.6 on journal page~109, in exactly the parameterisation used
below: ``Characterize graphs minimizing/maximizing the EUSO index among all
fixed-order connected graphs with a given (i) diameter or (ii) domination
number.''

Throughout, \emph{given diameter} means diameter \emph{exactly} $D$: for
$n\ge2$ and $D\ge1$ put
\[
 \Gnd_{n,D}=\{G:\ G\ \text{connected},\ |V(G)|=n,\ \diam(G)=D\},
\]
\[
 M(n,D)=\max_{G\in\Gnd_{n,D}}\EU(G),
 \qquad
 m(n,D)=\min_{G\in\Gnd_{n,D}}\EU(G).
\]

\section{Results}

For $a,b\ge1$ with $a+b=n-2$ let $H(a,b)$ be the graph on
$\{u,v\}\cup W$, $|W|=n-2$, obtained from the complete graph on $W$ by
splitting $W=A\sqcup B$ with $|A|=a$, $|B|=b$, joining $u$ to every vertex of
$A$ and $v$ to every vertex of $B$, and leaving $u$ and $v$
\emph{non-adjacent}. Note $H(a,b)\cong H(b,a)$.

\begin{theorem}\label{thm:B}
For every $n\ge4$,
\[
 M(n,3)=\frac{\sqrt3\,(n-2)^2(n-3)}{2}
        +\sqrt{n^2-3n+3}
        +(n-3)\sqrt{3n^2-15n+19},
\]
and the maximum is attained by exactly one graph up to isomorphism, namely
$H(1,n-3)$.
\end{theorem}

\begin{theorem}\label{thm:A}
$m(9,3)=8\sqrt3+6\sqrt{19}=40.009800121795059\ldots$, attained by exactly one
graph up to isomorphism, the bicyclic theta graph $\Theta(3,3,4)$ of
\S\ref{sec:objects}. In particular the minimiser of $\Gnd_{9,3}$ is neither a
tree nor unicyclic: the best tree in the cell is the double star $D(3,4)$ with
$\EU=\sqrt{61}+3\sqrt{21}+4\sqrt{31}=43.8290342121\ldots$, only third, behind
the unicyclic runner-up $4\sqrt{21}+8\sqrt3+4\sqrt7=42.7697144846\ldots$.
\end{theorem}

At $n=9$ Theorem~\ref{thm:B} reads
\[
 M(9,3)=147\sqrt3+\sqrt{57}+6\sqrt{127}=329.777869165403581\ldots,
\]
attained uniquely up to isomorphism by $H(1,6)$, with margin
$6.5941623645\ldots$ over the runner-up
$\EU(H(2,5))=147\sqrt3+2\sqrt{67}+5\sqrt{109}$.

Theorem~\ref{thm:B} is proved in \S\ref{sec:max} by a closed-form argument that
uses nothing about $n$ beyond $n\ge4$ and performs no search. The uniqueness
and exhaustiveness assertions of Theorem~\ref{thm:A} are computer-assisted: the
\emph{value} $8\sqrt3+6\sqrt{19}$ and the object $\Theta(3,3,4)$ are
hand-checkable from \S\ref{sec:objects}, and \S\ref{sec:min} confines the
minimiser to $m\le10$ edges by a strict inequality; but ``no other member of
the cell does better, and exactly one attains the value'' rests on the
exhaustive census of \S\ref{sec:min}, shipped with this paper as
\texttt{census.c}, \texttt{census.py} and \texttt{census.output.txt}.

\medskip
\noindent\textbf{Status.} Problem~1 quantifies
over all $n$ and all $D$ and is \emph{not} resolved here. What is settled is
the maximum half of the column $D=3$ for all $n\ge4$, together with the single
cell $(n,D)=(9,3)$ in both directions. The minimum half at $D=3$ is undecided
already for $n=10$; no cell with $D\ge4$ and $n>9$ is touched. Every uniqueness
statement is up to isomorphism. See \S\ref{sec:open}.

\medskip
\noindent\textbf{Prior art, and what it does not settle.} We make no priority
claim. The extremal \emph{object} $H(1,n-3)$ is already in print for this
class: Vetr\'ik \cite{Vetrik} proves sharp upper bounds for
$I_f(G)=\sum_{uv\in E}f(d(u),d(v))$ over graphs of given order and diameter,
and his corollary's maximiser for the sum-connectivity index $\chi_a$ at $a=1$
must be $H(1,n-3)$ as well, since $\chi_1$ is the first Zagreb index
$M_1=\sum_vd_v^2$, at $n=9$ $M_1(H(a,b))=343+a^2+b^2$ gives $380>372>368$ on
the three splits, and the whole tier $m\le27$ obeys
$M_1\le\Delta\cdot2m\le14m\le378<380$. What \cite{Vetrik} does not settle, so
far as we can tell, is $\EU$ itself: every index family named in its abstract
carries an exponent $\ge1$, whereas the edge function of \eqref{eq:EU} is
$1$-homogeneous and concave, and the Euler--Sombor index did not exist in 2023
\cite{TangEtAl}. That reading is an inference from an abstract and from the
restatement of the hypothesis in \cite{VetrikDML}, not a reading of the
theorem; \S\ref{sec:open} records the risk that it is wrong. The adjacent
Euler--Sombor items are these. Sekar, Balachandran, Elumalai and Liu
\cite{Sekar} bound $\EU$ over the \emph{unicyclic} graphs of given diameter, a
proper subclass of $\Gnd_{n,D}$, and so do not decide either extremum over the
whole cell; their Theorem~3.2 is used below only as a normalisation pin.
Kizilirmak \cite{Kizilirmak} works with $\EU$ under a diameter constraint;
only Lemma~2 of that paper, the value $\EU(C_9)$, is used here, again as a pin.
Khanra and Das \cite{KhanraDas} and the survey \cite{Survey} pose the problem
and do not answer it. So far as we could determine, none of these items
exhibits $M(n,3)$ or the minimiser of $\Gnd_{9,3}$; we did not read
\cite{Vetrik}.
What \S\ref{sec:max} supplies is the Euler--Sombor \emph{value} of $M(n,3)$
with an argument that the maximiser is unique, and \S\ref{sec:min} the minimum
in the single cell $(9,3)$. The auxiliary bound of Step~1 below is classical
and is credited there.

\section{The two extremal objects, exhibited}\label{sec:objects}

Both are printed completely, so a reader needs no code and no graph6 decoder.

\subsection*{The maximiser $H(1,6)$: $9$ vertices, $28$ edges}
Vertices $u,v,w_1,\ldots,w_7$. Edges: the $21$ pairs $w_iw_j$,
$1\le i<j\le7$; the edge $uw_1$; and the six edges $vw_2,vw_3,\ldots,vw_7$.
The vertices $u$ and $v$ are not adjacent.

Degrees: $d(u)=1$, $d(v)=6$, and $d(w_i)=7$ for $i=1,\ldots,7$; the degree
sequence is $(1,6,7,7,7,7,7,7,7)$. Checks: $21+1+6=28=\binom82$ and
$1+6+7\cdot7=56=2\cdot28$.

Diameter exactly $3$: the only neighbour of $u$ is $w_1$, and $w_1\not\sim v$,
so $d(u,v)=3$ along $u\,w_1\,w_j\,v$ for any $j\ge2$; every other pair of
vertices is at distance at most $2$, since $w_1,\ldots,w_7$ is a clique meeting
every neighbourhood.

Its index: each of the $21$ clique edges joins two vertices of degree $7$ and
contributes $\sqrt{49+49+49}=7\sqrt3$; the edge $uw_1$ joins degrees $(1,7)$
and contributes $\sqrt{1+49+7}=\sqrt{57}$; each edge $vw_i$ joins degrees
$(6,7)$ and contributes $\sqrt{36+49+42}=\sqrt{127}$. Hence
\begin{align*}
 \EU(H(1,6))&=147\sqrt3+\sqrt{57}+6\sqrt{127}\\
 &=254.6114687126+7.5498344353+67.6165660175\\
 &=329.7778691654\ldots
\end{align*}
The edge degree-pair profile is $(1,7)^1,(6,7)^6,(7,7)^{21}$.

\subsection*{The minimiser $\Theta(3,3,4)$: $9$ vertices, $10$ edges, bicyclic}
Two vertices $u,v$ joined by three internally disjoint paths of lengths
$3,3,4$:
\[
 u\,a_1\,a_2\,v,\qquad u\,b_1\,b_2\,v,\qquad u\,c_1\,c_2\,c_3\,v .
\]
Vertices: $u,v,a_1,a_2,b_1,b_2,c_1,c_2,c_3$. Edges ($10$):
$ua_1$, $a_1a_2$, $a_2v$, $ub_1$, $b_1b_2$, $b_2v$, $uc_1$, $c_1c_2$,
$c_2c_3$, $c_3v$.
Degrees: $d(u)=d(v)=3$ and all seven internal vertices have degree $2$; the
degree sum is $7\cdot2+2\cdot3=20=2\cdot10$, and the cyclomatic number is
$10-9+1=2$, so the graph is bicyclic.

Diameter exactly $3$: $u\not\sim v$ and the shortest $u$--$v$ path has length
$3$; the extremal pairs are two internal vertices on different paths, e.g.
$d(a_1,b_2)=3$ (via $a_1ub_1b_2$) and $d(a_1,c_2)=3$ (via $a_1uc_1c_2$), and no
pair exceeds $3$.

Its index: the four path-interior edges $a_1a_2,b_1b_2,c_1c_2,c_2c_3$ join two
degree-$2$ vertices and contribute $\sqrt{4+4+4}=2\sqrt3$ each; the six
remaining edges join a degree-$2$ vertex to $u$ or $v$ and contribute
$\sqrt{4+9+6}=\sqrt{19}$ each. Hence
\begin{align*}
 \EU(\Theta(3,3,4))&=8\sqrt3+6\sqrt{19}\\
 &=13.8564064606+26.1533936612=40.0098001218\ldots
\end{align*}
The profile is $(2,2)^4,(2,3)^6$.

\begin{remark}\label{rem:ties}
All four theta graphs $\Theta(l_1,l_2,l_3)$ on $9$ vertices with
$l_1+l_2+l_3=10$ and every $l_i\ge2$ share this profile and hence the
\emph{same} index $8\sqrt3+6\sqrt{19}$. What separates them is the diameter:
$\Theta(2,2,6)$, $\Theta(2,3,5)$ and $\Theta(2,4,4)$ all have diameter $4$,
and only $\Theta(3,3,4)$ has diameter $3$. So the minimiser is unique
\emph{in the cell} although its value is shared by graphs outside it; the
clause ``diameter exactly $3$'' does real work.
\end{remark}

\section{Proof of Theorem~\ref{thm:B}}\label{sec:max}

Fix $n\ge4$ and put $c=n-2$. In $H(a,b)$ we have $d(u)=a$, $d(v)=b$, every
vertex of $W$ has degree $c=n-2$, the size is
$\binom{n-2}2+a+b=\binom{n-1}2$, and $\diam H(a,b)=3$ exactly (the argument of
\S\ref{sec:objects} for $(a,b)=(1,6)$ is verbatim the general one).

\medskip
\noindent\textbf{Step 1 (edge ceiling; classical).}
Let $G\in\Gnd_{n,3}$ and pick $u,v$ with $d(u,v)=3$. Put $A=N(u)$, $B=N(v)$
and $R=V(G)\setminus(\{u,v\}\cup A\cup B)$. Since $d(u,v)=3$ we have
$u\not\sim v$ and $A\cap B=\emptyset$, so $|A|+|B|+|R|=n-2$. The pairs
$uv$; $\{u,b\}$ for $b\in B$; $\{u,r\}$ and $\{v,r\}$ for $r\in R$; and
$\{v,a\}$ for $a\in A$ are pairwise distinct non-edges, and there are
$1+|A|+|B|+2|R|=(n-1)+|R|\ge n-1$ of them. Hence
\begin{equation}\label{eq:ore}
 |E(G)|\ \le\ \binom n2-(n-1)\ =\ \binom{n-1}2 .
\end{equation}

\medskip
\noindent\textbf{Step 2 (equality cases).}
Equality in \eqref{eq:ore} forces $R=\emptyset$ and forces every pair not on
the list above to be an edge, i.e.\ $A\cup B$ induces a complete graph, $u$ is
joined to all of $A$ and $v$ to all of $B$. Thus $G=H(|A|,|B|)$ with
$|A|+|B|=n-2$ and $|A|,|B|\ge1$ (if $|A|=0$ then $u$ is isolated). Up to
isomorphism these are $\lfloor(n-2)/2\rfloor$ graphs; at $n=9$ exactly three,
$H(1,6)$, $H(2,5)$, $H(3,4)$.

\begin{remark}[Attribution]\label{rem:ore}
Steps 1 and 2 --- the bound \eqref{eq:ore} for a graph of diameter $3$
\emph{and} its equality characterisation --- are classical: they are the
Ore-type extremal problem ``maximum size of a graph of given order and
diameter''. They are re-derived here because the derivation is four lines, not
because they are new; see \cite{Ore2023} for a modern proof and entry point,
and \cite{Tokyo2010} for the $2$-connected refinement. What is proved below is
the Euler--Sombor layer on top: Step~3 and Step~4.
\end{remark}

\medskip
\noindent\textbf{Step 3 (the whole tier $|E|\le\binom{n-1}2-1$ is beaten).}
If $\diam(G)\ge3$ then no vertex has degree $n-1$ --- such a vertex is adjacent
to everything and forces $\diam\le2$ --- so $\Delta(G)\le n-2$ and every edge
term of \eqref{eq:EU} obeys
$\sqrt{d_x^2+d_xd_y+d_y^2}\le\sqrt{3\Delta^2}=\sqrt3\,(n-2)$.
(The hypothesis is $\diam\ge3$, not $\diam\ge2$: the star $K_{1,n-1}$ has
diameter $2$ and $\Delta=n-1$.)
Since $\binom{n-1}2-1-\binom{n-2}2=n-3$, any $G\in\Gnd_{n,3}$ with
$|E(G)|\le\binom{n-1}2-1$ satisfies
\[
 \EU(G)\ \le\ \underbrace{\sqrt3(n-2)\binom{n-2}2}_{\text{the clique part of every }H(a,b)}
            +\ \sqrt3\,(n-2)(n-3),
\]
so that whole tier is strictly beaten by $H(1,n-3)$ as soon as
\begin{equation}\label{eq:star}
\begin{split}
 \sqrt3\,(n-2)(n-3)\ <\ &\sqrt{(n-2)^2+(n-2)+1}\\
 &+(n-3)\sqrt{3(n-3)^2+3(n-3)+1}.
\end{split}
\end{equation}
Write $s=n-3\ge1$; \eqref{eq:star} becomes
$\sqrt{s^2+3s+3}+s\sqrt{3s^2+3s+1}>\sqrt3\,s(s+1)$. Since
$\sqrt3\,s(s+1)=s\sqrt{3s^2+6s+3}$,
\begin{align*}
 \sqrt3\,s(s+1)-s\sqrt{3s^2+3s+1}
 &=\frac{s(3s+2)}{\sqrt{3s^2+6s+3}+\sqrt{3s^2+3s+1}}\\
 &<\frac{s(3s+2)}{2\sqrt{3s^2}}
 =\frac{\sqrt3}{2}s+\frac1{\sqrt3},
\end{align*}
while $\sqrt{s^2+3s+3}>s+1$ because $s^2+3s+3-(s+1)^2=s+2>0$. Finally
\begin{align*}
 (s+1)-\Bigl(\tfrac{\sqrt3}2s+\tfrac1{\sqrt3}\Bigr)
 &=\Bigl(1-\tfrac{\sqrt3}2\Bigr)s+\Bigl(1-\tfrac1{\sqrt3}\Bigr)\\
 &=0.1339745\ldots s+0.4226497\ldots>0 ,
\end{align*}
so \eqref{eq:star} holds for every $s\ge1$, with a margin growing linearly in
$n$. The left-hand bound is valid on the \emph{whole} discarded tier, so
nothing can hide there. At $n=9$ the cap on the tier is
$27\cdot7\sqrt3=189\sqrt3=327.3576026305\ldots$, below
$329.7778691654\ldots$

\medskip
\noindent\textbf{Step 4 (which $H(a,b)$ wins).}
Every edge from the apex $u$ of degree $t$ to the clique joins degrees $t$ and
$c$, so those $t$ edges contribute
\[
 \varphi(t)=t\sqrt{t^2+ct+c^2},
 \qquad\text{and}\qquad
 \EU(H(a,b))=\sqrt3\,c\binom c2+\varphi(a)+\varphi(c-a).
\]
With $q(t)=t^2+ct+c^2$,
\[
 \varphi'(t)=\frac{4t^2+3ct+2c^2}{2\sqrt q},
 \qquad
 \varphi''(t)=\frac{8t^3+12ct^2+15c^2t+4c^3}{4q^{3/2}}>0
\]
for $t\ge0$ and $c>0$,
the numerator of $\varphi''$ being
$2(8t+3c)q-(4t^2+3ct+2c^2)(2t+c)
 =(16t^3+22ct^2+22c^2t+6c^3)-(8t^3+10ct^2+7c^2t+2c^3)$.
So $\varphi$ is strictly convex, hence $a\mapsto\varphi(a)+\varphi(c-a)$ is
strictly convex on $[1,c-1]$ and is maximised only at the endpoints $a=1$ and
$a=c-1$, which give isomorphic graphs. Therefore $H(1,n-3)$ is the unique
maximiser among the graphs of Step~2, and by Step~3 among all of
$\Gnd_{n,3}$. Evaluating $\EU(H(1,n-3))$ gives
\begin{align*}
 \sqrt3\,(n-2)\binom{n-2}2&+\sqrt{1+(n-2)+(n-2)^2}\\
 &+(n-3)\sqrt{(n-3)^2+(n-3)(n-2)+(n-2)^2},
\end{align*}
which is the closed form of Theorem~\ref{thm:B}. \qed

\section{The minimum at $n=9$, $D=3$}\label{sec:min}

\noindent\textbf{No monotonicity is available.} Adding an edge strictly
increases $\EU$, but $\Gnd_{n,D}$ is closed under neither edge addition (an
added edge can drop the diameter) nor edge deletion (a deleted edge can raise
it or disconnect), so the minimiser is not forced onto a tree: $C_6$ has
diameter $3$ and every spanning tree of $C_6$ is $P_6$, of diameter $5$.

\medskip
\noindent\textbf{A floor that amputates the search space.}
For $x,y>0$ we have $x^2+xy+y^2-\tfrac34(x+y)^2=\tfrac14(x-y)^2\ge0$, so each
edge term of \eqref{eq:EU} is at least $\tfrac{\sqrt3}2(d_x+d_y)$; summing over
edges and applying Cauchy--Schwarz to $\sum_vd_v^2\ge4m^2/n$ gives, for every
graph of order $n$ and size $m$,
\begin{equation}\label{eq:floor}
 \EU(G)\ \ge\ \frac{\sqrt3}{2}\sum_{v}d_v^2\ \ge\ \frac{2\sqrt3\,m^2}{n}.
\end{equation}
At $n=9$, $m=11$ already gives
\[
 \frac{2\sqrt3\cdot121}{9}=46.5729217146\ldots>43.8290342121\ldots=\EU(D(3,4)),
\]
and the right-hand side of \eqref{eq:floor} increases in $m$. Since $D(3,4)$ lies
in $\Gnd_{9,3}$, the minimiser has $m\in\{8,9,10\}$: it is a tree, unicyclic
or bicyclic.

\medskip
\noindent\textbf{$m=8$ by hand.} The only $9$-vertex trees of diameter $3$ are
the double stars $D(a,7-a)$ (adjacent centres of degrees $a+1$ and $8-a$,
carrying $a$ and $7-a$ leaves), and
\[
 \begin{array}{l}
 D(3,4)=\sqrt{61}+3\sqrt{21}+4\sqrt{31}=43.8290342121\ldots,\\
 D(2,5)=\sqrt{63}+2\sqrt{13}+5\sqrt{43}=47.9355491\ldots,\\
 D(1,6)=\sqrt{67}+\sqrt7+6\sqrt{57}=56.1301107\ldots,
 \end{array}
\]
so the best tree in the cell is $D(3,4)$, already above
$8\sqrt3+6\sqrt{19}$.

\medskip
\noindent\textbf{$m=10$: the cheapest profile is not realisable at diameter
$3$.}
By \eqref{eq:floor} the cheapest way to spend $10$ edges on $9$ vertices is to
keep the degrees as equal as the class allows, and the cheapest degree sequence
is $(3,3,2^7)$. Among edge profiles realising it, the cheapest is the one with
a single $(3,3)$-edge, worth
$13\sqrt3+4\sqrt{19}=39.9522\ldots<40.0098\ldots$; so if that profile were
realisable at diameter $3$ the claimed minimum would be wrong. It is not
realisable. A single $(3,3)$-edge makes the two degree-$3$ vertices adjacent,
which forces either a theta graph $\Theta(1,b,c)$ with $b+c=9$ or a dumbbell
with a handle of length $1$; and $\Theta(1,2,7)$, $\Theta(1,3,6)$,
$\Theta(1,4,5)$ and every $C_a$--$C_b$ dumbbell with $a+b=9$ have diameter at
least $4$. If, on the other hand, the profile has no
$(3,3)$-edge, then the two degree-$3$ vertices are non-adjacent, all six of
their edges run to degree-$2$ vertices, and a connected graph with degrees in
$\{2,3\}$ and exactly two vertices of degree $3$ is a theta graph or a
dumbbell; a $9$-vertex dumbbell has diameter at least $5$, so the graph is
$\Theta(l_1,l_2,l_3)$ with $\sum l_i=10$ and every $l_i\ge2$, and by
Remark~\ref{rem:ties} only $(3,3,4)$ has diameter $3$.
No written argument is given here for the remaining degree sequences with
$m=10$ --- the next one met is $(4,2^8)$, whose cheapest profile costs
$4\sqrt{28}+12\sqrt3=41.9506\ldots$, and we do not go through the rest --- and
none is given for $m=9$. Those cases are settled by the census below, not by
hand.

\medskip
\noindent\textbf{$m=9$, exhaustiveness and uniqueness: the census.} The
unicyclic case and the exhaustiveness and uniqueness assertions of
Theorem~\ref{thm:A} rest on enumeration. That enumeration is shipped with this
paper: \texttt{census.c} is the enumeration kernel, \texttt{census.py} the
driver carrying the arithmetic, the controls, the automorphism counts and the
verdict, and \texttt{census.output.txt} the retained output of the run, whose
header records the reproduction command. It has two passes that share no
enumeration machinery. Pass~1 reads every isomorphism class of connected
nine-vertex graph from \texttt{nauty}'s \texttt{geng}, stratifies them by
diameter, and compares $\EU$ over the cell in $80$-significant-digit
\texttt{Decimal} arithmetic; since \texttt{geng} emits one graph6 line per
class, the number of classes in an extremal bucket \emph{is} the number of
ties, and it is $1$ at each end. Pass~2 uses no generator, no canonical form
and no isomorphism test: it sweeps all $2^{36}$ labelled graphs on nine
vertices, finds no graph of the cell beating either extremum, and finds exactly
$504=9!/|\mathrm{Aut}(H(1,6))|=9!/720$ labelled graphs within $10^{-6}$ of the
maximum and exactly $90\,720=9!/|\mathrm{Aut}(\Theta(3,3,4))|=9!/4$ within
$10^{-6}$ of the minimum, i.e.\ one orbit at each end and nothing else. All
$43$ checks pass.

The driver reproduces its controls before printing any extremum and refuses to
print one if a control fails. \texttt{geng -c 9} emits
$261\,080=A001349(9)$ connected classes; the eight diameter strata are
$1$, $91\,518$, $148\,229$, $19\,320$, $1818$, $180$, $13$, $1$, which sum to
that figure, so the cell
$\Gnd_{9,3}$ holds exactly $148\,229$ isomorphism classes. The $21$ per-size
class counts of the $D=3$ stratum, for $m=8,\ldots,28$,
\[
 \begin{gathered}
 3,\ 17,\ 76,\ 369,\ 1483,\ 4381,\ 9621,\ 16547,\ 23013,\ 26173,\ 24441,\\
 18828,\ 12040,\ 6477,\ 2980,\ 1185,\ 412,\ 131,\ 38,\ 11,\ 3,
 \end{gathered}
\]
sum to $148\,229$ as well; and the labelled pass's own count of connected
labelled graphs on nine nodes, $66\,296\,291\,072$, matches the exponential
formula evaluated in exact integer arithmetic. Two published values by other
authors pin the normalisation \eqref{eq:EU}: the census's maximum over the
\emph{unicyclic} members of $\Gnd_{9,3}$ is the right-hand side of Theorem~3.2
of \cite{Sekar} at $n=9$, namely
$5\sqrt{57}+\sqrt{79}+\sqrt{67}+\sqrt{13}+\sqrt{19}=62.7871695845\ldots$, so
that theorem is sharp at $n=9$; and the census's minimum over $\Gnd_{9,4}$ is
Lemma~2 of \cite{Kizilirmak}, $\EU(C_9)=18\sqrt3=31.1769145362\ldots$,
attained by the $9$-cycle alone. The census also reports that the smallest gap
between two distinct profile values anywhere in the cell is
$1.0380112\ldots\times10^{-7}$; at $80$ significant digits no comparison it
makes is delicate.

\begin{remark}
We claim the minimiser of no cell other than $(9,3)$.
\end{remark}

\section{What is not settled}\label{sec:open}

\noindent\textbf{The one risk to Theorem~\ref{thm:B}.} Vetr\'ik \cite{Vetrik}
proves sharp \emph{upper} bounds on $I_f$ over graphs of given order and
diameter, and \S2 records that his extremal object must coincide with ours. Its
full text is paywalled, with no open-access or repository copy and no preprint,
and we could not read it. What we did read is the publisher's abstract, in which
every index family carries an exponent $\ge1$ ($\chi_a$, $R_a$, $FGO_a$ with
$a\ge1$; $GZ_{a,b}$, $M_{a,b}$ with $a,b\ge1$), whereas the edge function of
\eqref{eq:EU} is $1$-homogeneous and concave. Vetr\'ik's own later paper
\cite{VetrikJAMC} says of \cite{Vetrik} that ``the result presented in [17]
holds for $R_a$, where $a\ge1$, and for $M_{a,b}$, where $a,b\ge1$'', and the
$f$-condition as printed in \cite{VetrikDML} fails for
$f(x,y)=\sqrt{x^2+xy+y^2}$: with $(x_1,y_1)=(1,1)$, $(x_2,y_2)=(1,5)$ one has
$\sqrt7-\sqrt3=0.9137\ldots>\sqrt{39}-\sqrt{31}=0.6772\ldots$, so the required
increasing-differences clause is violated. On that evidence $\EU$ appears to lie
outside the class. \emph{But that is an inference from an abstract and from a
second paper's restatement of the hypothesis, not a reading of the theorem.} If
the hypothesis of \cite{Vetrik} does admit a concave $1$-homogeneous $f$, then
Theorem~\ref{thm:B} is a rediscovery and what survives is
Theorem~\ref{thm:A}'s minimum half together with the two numerical values.

\medskip
\noindent\textbf{Cells left open.} The minimum half at $D=3$ is undecided for
every $n\ge10$, and every cell with $D\ge4$ and $n>9$ is untouched in both
directions.

\medskip
\noindent\textbf{What the two shipped programs do and do not cover.} The census
of \S\ref{sec:min} --- \texttt{census.c}, \texttt{census.py} and its retained
transcript \texttt{census.output.txt}, whose header records the sha256 sums of
the two files that were run, the wall clock and the reproduction
command --- is what carries the exhaustiveness and the uniqueness of both
extrema of $\Gnd_{9,3}$. The second program, \texttt{verify.py} with its
transcript \texttt{verify.output.txt}, enumerates nothing: it re-derives what a
pencil can reach --- Theorem~\ref{thm:B} in full, the two objects of
\S\ref{sec:objects} with their closed-form values and diameters, the strict
ordering $40.0098\ldots<42.7697\ldots<43.8290\ldots$ among the three named
graphs $\Theta(3,3,4)$, one hand-built unicyclic member of the cell and the
double star $D(3,4)$, and no others, the amputation \eqref{eq:floor}, the tier
cap of Step~3 and the strata identities
--- and its closing note says correctly that it cannot itself see the
exhaustiveness or the uniqueness of the minimum, and that the strata counts it
re-prints are transcriptions of the census rather than re-reads. Neither program
reads \cite{Vetrik}; the $f$-class exclusion discussed above is a computation on
the condition as printed in \cite{VetrikDML}, not a reading of Vetr\'ik's
hypothesis. Where a
surd identity is asserted it is checked to a stated tolerance.

\begin{thebibliography}{9}

\bibitem{Survey}
A.~Ali, I.~Gutman, T.~R\'eti, A.~M. Albalahi, and A.~E. Hamza,
\emph{Survey on extremal results and bounds for elliptic Sombor and
Euler--Sombor indices},
Discrete Math. Lett. \textbf{17} (2026), 101--112,
\href{https://doi.org/10.47443/dml.2026.056}{doi:10.47443/dml.2026.056}.
Problem~3.1.6 is on journal page~109.

\bibitem{KhanraDas}
B.~Khanra and S.~Das,
\emph{Euler Sombor index of trees, unicyclic and chemical graphs},
MATCH Commun. Math. Comput. Chem. \textbf{94} (2025), 525--548.
The open problem quoted in \S1 is item~(1) of its concluding list, journal
page~546.

\bibitem{Kizilirmak}
G.~O. Kizilirmak,
MATCH Commun. Math. Comput. Chem. \textbf{94} (2025), 725--738.
Its abstract, journal page~725, prints the same normalisation of $\EU$ as
\eqref{eq:EU}; the Lemma~2 used as a pin in \S\ref{sec:min} is on journal
page~728. Title not recorded in the source we consulted.

\bibitem{Ore2023}
\emph{A simple proof of Ore's theorem on the maximum size of $k$-connected
graphs with given order and diameter},
Graphs Combin. (2023),
\href{https://doi.org/10.1007/s00373-023-02628-w}{doi:10.1007/s00373-023-02628-w}.
Authors not recorded in the source we consulted.

\bibitem{Sekar}
S.~Sekar, S.~Balachandran, S.~Elumalai, and H.~Liu,
\emph{Maximum Euler Sombor index of unicyclic graphs with given diameter},
MATCH Commun. Math. Comput. Chem. \textbf{95} (2026), no.~3, 759--776,
\href{https://doi.org/10.46793/match.95-3.12225}{doi:10.46793/match.95-3.12225}.
Problem~1 is on journal page~762; the index \eqref{eq:EU} is defined on
journal page~761; Theorem~3.2 is on journal page~764.

\bibitem{TangEtAl}
Z.~Tang, Q.~Li, and H.~Deng,
Int. J. Quantum Chem. \textbf{124} (2024), e27387; and I.~Gutman,
Milit. Techn. Courier \textbf{72} (2024), 1--12.
These are the two 2024 papers that introduce the Euler--Sombor index, cited as
references [20] and [7] of \cite{Sekar}. Titles not recorded in the source we
consulted.

\bibitem{Tokyo2010}
\emph{Extremal $2$-connected graphs with given diameter},
Tokyo J. Math. (2010),
\href{https://doi.org/10.3836/tjm/1270041609}{doi:10.3836/tjm/1270041609}.
Authors not recorded in the source we consulted.

\bibitem{Vetrik}
T.~Vetr\'ik,
\emph{Degree-based function index for graphs with given diameter},
Discrete Appl. Math. \textbf{333} (2023), 59--70,
\href{https://doi.org/10.1016/j.dam.2023.02.018}{doi:10.1016/j.dam.2023.02.018}.
Nearest prior art; see \S2 and \S\ref{sec:open}. Its full text was not
accessible to us.

\bibitem{VetrikDML}
T.~Vetr\'ik and P.~Kaemawichanurat,
Discrete Math. Lett. \textbf{15} (2025), 23--30.
The $f$-condition tested in \S\ref{sec:open} is Definition~1.1 as printed
here, and its Lemma~1.1 attributes three cases to \cite{Vetrik}. Title not
recorded in the source we consulted.

\bibitem{VetrikJAMC}
T.~Vetr\'ik,
J. Appl. Math. Comput. (2024),
\href{https://doi.org/10.1007/s12190-024-02307-w}{doi:10.1007/s12190-024-02307-w};
open access, and the source of the sentence about \cite{Vetrik} quoted in
\S\ref{sec:open}. Title not recorded in the source we consulted.

\end{thebibliography}

\end{document}
